% Einheit 2 von 5 – Vektoren unter Matrizen (bank/matrizen-und-uebergangsprozesse/e2.jsonl, zusammenbau v0.1)
%% TODO Zeitmarke Einheit 2: Marken-Zeile nicht lesbar
%% TODO Zweigzeile Teil 1: was hier gelernt wird (Satz in Schülersprache aus dem Einheitstitel) – Einheit 2
\einheitenkopf[e2]{Einheit 2 von 5 · Vektoren unter Matrizen}
\zweigzeile{keine Prüfungsaufgabe}
\setcounter{aufgabe}{10}

%% TODO Titel: Ich-kann-Satz fehlt in der Bank (Platzhalter: Kettenname) – Nr. 11
%% TODO Anweisung: ein Satz über der ersten Teilaufgabe fehlt in der Bank – Nr. 11
\begin{aufgabe}{Vektoren unter Matrizen}
%% TODO \rechenplatz{n}: Schrittzahl des Grundfalls fehlt in der Bank (nur bei fester Schreibform, 2.3 b)
\begin{teile}
\teil Allgemein oder am Beispiel? Nur ankreuzen: „Zeige, dass jede Matrix $((p | q), (r | s))$ mit $p + r = 1$ und $q + s = 1$ die Komponentensumme jedes Vektors erhält.“ \\ \kreuz{allgemein zeigen} \\ \kreuz{ein Beispiel genügt}
\end{teile}
\begin{gleichungsraster}[2]
\gl{\text{Bestimme alle Vektoren v mit $M \cdot v = v$ für $M = ((0{,}6 | 0{,}2), (0{,}4 | 0{,}8))$.}} &
\gl{\text{Bestimme alle Vektoren v mit $M \cdot v = 3 \cdot v$ für $M = ((2 | 1), (0 | 3))$.}} \\
\gl{\text{$M = ((1 | 2 | 0), (0 | 1 | 1), (2 | 0 | 1))$ und $M \cdot (a | 1 | 2) = (5 | 3 | 8)$. Bestimme a und prüfe mit den übrigen Zeilen.}} &
\end{gleichungsraster}
\begin{teile}
\teil Zeige allgemein: Ist $M = ((p | q), (r | s))$ mit $p + r = 1$ und $q + s = 1$, so hat $M \cdot (x | y)$ dieselbe Komponentensumme wie $(x | y)$.
\end{teile}
\begin{gleichungsraster}[2]
\gl{\text{Eine Matrix heißt selbstinvers, wenn $M \cdot M = E$ gilt. Für welche a ist $M = ((a | 1), (0 | -1))$ selbstinvers?}} &
\end{gleichungsraster}
\begin{teile}
\teil $M = ((1 | 2), (2 | 4))$. Zeige, dass $M \cdot v = (0 | 0)$ außer dem Nullvektor weitere Lösungen hat, und gib alle an.
\teil $A = ((3 | 2), (1 | 4))$ und $w = (1 | 1)$ mit $A \cdot w = 5 \cdot w$. Für jede Zahl c gibt es eine Zahl b mit $A^{-1} \cdot (b \cdot w) = c \cdot w$. Bestimme b in Abhängigkeit von c, ohne $A^{-1}$ zu berechnen \hfill (Abitur 2026 LK).
\end{teile}
\end{aufgabe}

%% TODO Titel: Ich-kann-Satz fehlt in der Bank (Platzhalter: Kettenname) – Nr. 12
%% TODO Anweisung: ein Satz über der ersten Teilaufgabe fehlt in der Bank – Nr. 12
\begin{aufgabe}{Vektoren unter Matrizen – Fehler finden, Begründen}
\begin{teile}
\teil Jan soll alle v mit $M \cdot v = v$ für $M = ((0{,}5 | 0{,}25), (0{,}5 | 0{,}75))$ bestimmen. Er schreibt: „$v = (0 | 0)$, denn der Nullvektor erfüllt die Gleichung.“ Finde den Fehler und gib alle Lösungen an.
\teil Begründe: Erfüllt v die Gleichung $M \cdot v = t \cdot v$, dann erfüllt auch $5v$ sie.
\end{teile}
\end{aufgabe}
